\begin{explanationbox}
\textbf{\textcolor{CExplanation}{一句话直觉}}
余弦函数 $\cos x$ 的图像，在 $x=0$ 附近比对应的二次函数 $1-\frac{x^2}{2}$ 的图像“更鼓”或“更靠上”。

\medskip
\textbf{\textcolor{CExplanation}{核心拆解}}
\begin{description}[style=nextline, leftmargin=2.8em, labelsep=0.5em, itemsep=0.45em]
    \item[\textbf{要点一（比较对象）：}] 比较的对象是余弦函数 $\cos x$ 和二次函数 $y = 1 - \frac{x^2}{2}$。
    \item[\textbf{要点二（不等关系）：}] 对于所有实数 $x$，始终有 $\cos x$ 大于或等于 $1 - \frac{x^2}{2}$。
    \item[\textbf{要点三（相切条件）：}] 两个函数仅在 $x=0$ 时函数值相等，图像在该点相切。
\end{description}

\medskip
\textbf{\textcolor{CExplanation}{几何本质（切线放缩）}}
该不等式表明，二次函数 $y=1-\frac{x^2}{2}$ 是 $y=\cos x$ 在 $x=0$ 处的一条“下方切线”（更准确地说，是二阶泰勒多项式下界）。在 $x=0$ 附近，可以用这个简单的二次函数来近似或控制 $\cos x$ 的值，且保证是下界估计。

\medskip
\textbf{\textcolor{CExplanation}{代数意义（多项式逼近）}}
从代数上看， $1-\frac{x^2}{2}$ 是 $\cos x$ 的泰勒展开式（麦克劳林公式）的前两项。完整展开为 $\cos x = 1 - \frac{x^2}{2!} + \frac{x^4}{4!} - \cdots$，由于后续项正负交替且当 $x$ 较小时绝对值递减，前两项给出了一个全局下界。

\medskip
\textbf{\textcolor{CExplanation}{考点价值（放缩工具）}}
\begin{itemize}[leftmargin=*, itemsep=0.5em, topsep=0.4em]
    \item \textbf{考法一（证明不等式）：} 在证明涉及 $\cos x$ 的不等式时，可直接用它放缩，将超越式化为多项式。
    \item \textbf{考法二（估值或比较）：} 用于快速比较 $\cos x$ 与某个二次式的大小，或对 $\cos x$ 进行下界估计。
\end{itemize}

\medskip
\textbf{\textcolor{CExplanation}{顿悟点}}
\textbf{这个不等式将复杂的三角函数放缩为简单的二次函数，是“以直代曲”思想的一种体现，但这里“直”换成“抛物线”。它为我们处理 $\cos x$ 提供了一个普适、简洁的下界工具。}

\medskip
\textbf{\textcolor{CExplanation}{使用场景}}
\begin{itemize}[leftmargin=*, itemsep=0.5em, topsep=0.4em]
    \item \textbf{场景一（直接放缩）：} 当题目需要证明 $\cos x \ge g(x)$ 时，若 $g(x) \le 1-\frac{x^2}{2}$，则可利用此结论进行传递。
    \item \textbf{场景二（构建不等式链）：} 在连续放缩中，作为中间步骤，连接三角函数与多项式。
\end{itemize}
\end{explanationbox}
